\section*{Chapitre \ref{ch:b1:functions} --- Fonctions usuelles}

\begin{solution}{exo:b1:functions:1}
$\arcsin\frac{\sqrt 3}{2} = \frac\pi3$~;
$\;\arccos\bigl(-\frac12\bigr) = \frac{2\pi}{3}$~;
$\;\arctan(-1) = -\frac\pi4$.

$\arcsin\bigl(\sin\frac{5\pi}{6}\bigr)$~: $\sin\frac{5\pi}{6} =
\frac12$, et l'angle de $\intcc{-\frac\pi2}{\frac\pi2}$ dont le sinus
vaut $\frac12$ est $\frac\pi6$ (et non $\frac{5\pi}{6}$).

$\arccos\bigl(\cos\frac{7\pi}{4}\bigr)$~: $\cos\frac{7\pi}{4} =
\frac{\sqrt 2}{2}$, et l'angle de $\intcc{0}{\pi}$ ayant ce cosinus
est $\frac\pi4$.
\end{solution}

\begin{solution}{exo:b1:functions:2}
$f(x) = \arcsin(2x - 1)$ exige $-1 \leq 2x - 1 \leq 1$, c'est-à-dire
$x \in \intcc{0}{1}$. Sur $\intoo{0}{1}$, la dérivation des fonctions
composées et la \cref{prop:b1:functions:arcderiv} donnent
\[
f'(x) = \frac{2}{\sqrt{1 - (2x-1)^2}} = \frac{2}{\sqrt{4x - 4x^2}}
= \frac{1}{\sqrt{x(1 - x)}} .
\]
$g(x) = \arctan\sqrt x$ est définie pour $x \geq 0$, et pour $x > 0$~:
\[
g'(x) = \frac{1}{1 + x} \cdot \frac{1}{2\sqrt x}
= \frac{1}{2\sqrt x\,(1 + x)} .
\]
\end{solution}

\begin{solution}{exo:b1:functions:3}
$\cosh x + \sinh x = \frac{\eu^x + \eu^{-x}}{2} + \frac{\eu^x -
\eu^{-x}}{2} = \eu^x$. Donc, pour $n \in \Z$~:
\[
(\cosh x + \sinh x)^n = (\eu^x)^n = \eu^{nx}
= \cosh nx + \sinh nx .
\]
\end{solution}

\begin{solution}{exo:b1:functions:4}
Avec $y = \arcsin x$~: $\cos 2y = 1 - 2\sin^2 y = 1 - 2x^2$, donc
$\cos(2\arcsin x) = 1 - 2x^2$.

Avec $y = \arctan x$~: $\sin 2y = 2 \sin y \cos y = 2 \tan y \cos^2 y
= \frac{2\tan y}{1 + \tan^2 y}$, donc $\sin(2\arctan x) =
\dfrac{2x}{1 + x^2}$.
\end{solution}

\begin{solution}{exo:b1:functions:5}
De la plus lente à la plus rapide~:
\[
(\ln x)^{1000} \;\ll\; x^{100} \;\ll\; x^{\ln x} \;\ll\;
\eu^{\sqrt x} \;\ll\; \eu^{x/100} .
\]
Justifications~: $(\ln x)^{1000}/x^{100} \to 0$ d'après la
\cref{prop:b1:functions:powerrules} (les logarithmes perdent contre
les puissances). $x^{100} = \eu^{100\ln x}$ et $x^{\ln x} =
\eu^{(\ln x)^2}$~: comme $(\ln x)^2 - 100\ln x \to +\infty$, c'est la
seconde qui l'emporte. $x^{\ln x} = \eu^{(\ln x)^2} \ll \eu^{\sqrt x}$
car $(\ln x)^2/\sqrt x \to 0$ (les logarithmes perdent contre la
puissance $x^{1/4}$, le tout élevé au carré). Enfin $\sqrt x - x/100
\to -\infty$, donc $\eu^{\sqrt x} \ll \eu^{x/100}$.
\end{solution}

\begin{solution}{exo:b1:functions:6}
\omterm{def:b1:logic:sets}{Ensemble} de définition~: $x \neq \pm 1$, soit trois intervalles. Sur
chacun,
\[
f'(x) = \frac{\bigl(\frac{2x}{1-x^2}\bigr)'}{1 +
\frac{4x^2}{(1-x^2)^2}}
= \frac{\frac{2(1-x^2) + 4x^2}{(1-x^2)^2}}
{\frac{(1-x^2)^2 + 4x^2}{(1-x^2)^2}}
= \frac{2(1 + x^2)}{(1 + x^2)^2} = \frac{2}{1 + x^2}
= (2\arctan x)' .
\]
Donc $f(x) - 2\arctan x$ est constante sur chaque intervalle. Valeurs~:
en $x = 0$, $f(0) = 0$~: la constante vaut $0$ sur $\intoo{-1}{1}$.
Quand $x \to +\infty$, $\frac{2x}{1-x^2} \to 0^-$ donc $f \to 0$,
tandis que $2\arctan x \to \pi$~: la constante vaut $-\pi$ sur
$\intoo{1}{+\infty}$. Par imparité, elle vaut $+\pi$ sur
$\intoo{-\infty}{-1}$. En résumé~: $f = 2\arctan x$ sur
$\intoo{-1}{1}$, $f = 2\arctan x - \pi$ pour $x > 1$, $f = 2\arctan x
+ \pi$ pour $x < -1$. (C'est la formule de duplication de la tangente,
lue à travers $\arctan$.)
\end{solution}

\begin{solution}{exo:b1:functions:7}
$\cosh x = 2$~: avec $u = \eu^x > 0$, $u + u^{-1} = 4$, donc $u^2 - 4u
+ 1 = 0$, $u = 2 \pm \sqrt 3$~: $x = \ln(2 + \sqrt 3)$ ou $x = \ln(2 -
\sqrt 3) = -\ln(2 + \sqrt 3)$ (les deux solutions sont opposées,
puisque $\cosh$ est paire~; toutes deux conviennent). De façon
équivalente, $x = \pm \operatorname{arcosh} 2$.

$5\cosh x - 4\sinh x = 3$~: en substituant les définitions
exponentielles, $\frac{5(u + u^{-1}) - 4(u - u^{-1})}{2} = 3$,
c'est-à-dire $u + 9u^{-1} = 6$, c'est-à-dire $u^2 - 6u + 9 = (u -
3)^2 = 0$~: $u = 3$, une unique solution $x = \ln 3$.
\end{solution}

\begin{solution}{exo:b1:functions:8}
Posons $\alpha = \arctan\frac12 + \arctan\frac13$. Formule
d'addition~:
\[
\tan\alpha = \frac{\frac12 + \frac13}{1 - \frac12\cdot\frac13}
= \frac{5/6}{5/6} = 1 .
\]
Les deux \omterm{def:b1:functions:arc}{arctangentes} appartiennent à $\intoo{0}{\frac\pi4}$ (leurs
arguments sont dans $\intoo{0}{1}$), donc $\alpha \in
\intoo{0}{\frac\pi2}$~; le seul angle de cet intervalle dont la
tangente vaut $1$ est $\frac\pi4$.

Machin~: posons $\beta = \arctan\frac15$. Duplication, deux fois~:
\[
\tan 2\beta = \frac{2/5}{1 - 1/25} = \frac{5}{12},
\qquad
\tan 4\beta = \frac{2 \cdot 5/12}{1 - 25/144} = \frac{120}{119}.
\]
Puis
\[
\tan\Bigl(4\beta - \frac\pi4\Bigr)
= \frac{\frac{120}{119} - 1}{1 + \frac{120}{119}}
= \frac{1/119}{239/119} = \frac{1}{239}.
\]
Localisation~: $\beta < \arctan 1 = \frac\pi4$, et en fait
$\tan 4\beta = \frac{120}{119}$ est proche de $1$ avec $4\beta \in
\intoo{0}{\frac\pi2}$ (car $\beta < \frac\pi8$, puisque
$\tan\frac\pi8 = \sqrt 2 - 1 > \frac15$)~; donc $4\beta - \frac\pi4
\in \intoo{-\frac\pi4}{\frac\pi4}$, intervalle sur lequel $\arctan$
inverse $\tan$~: $4\beta - \frac\pi4 = \arctan\frac{1}{239}$, ce qui
est la \omterm{rem:b1:functions:interlude}{formule de Machin}.
\end{solution}

\begin{solution}{exo:b1:functions:9}
Posons $f(x) = x - \sin x$~: $f(0) = 0$ et $f'(x) = 1 - \cos x \geq
0$, donc $f \geq 0$ sur $\R_+$~: $\sin x \leq x$.

Posons $g(x) = \sin x - x + \frac{x^3}{6}$~: $g(0) = 0$, $g'(x) =
\cos x - 1 + \frac{x^2}{2}$, $g'(0) = 0$, $g''(x) = -\sin x + x =
f(x) \geq 0$ sur $\R_+$. Donc $g'$ est croissante avec $g'(0) = 0$,
d'où $g' \geq 0$, d'où $g$ croissante avec $g(0) = 0$~: $g \geq 0$,
c'est-à-dire $x - \frac{x^3}{6} \leq \sin x$ sur $\R_+$.

Par conséquent $0 \leq \frac{x - \sin x}{x^3} \leq \frac 16$ pour
$x > 0$~: le rapport est piégé à l'échelle de $\frac16$, et le
\cref{ch:b1:taylor} montre que sa limite vaut exactement $\frac 16$.
\end{solution}

\begin{solution}{exo:b1:functions:10}
Écrivons $y = \arcsin x$, de sorte que $2x\sqrt{1 - x^2} = 2\sin y\cos
y = \sin 2y$ et $h(x) = \arcsin(\sin 2y) - 2y$.

Pour $\abs x \leq \frac{\sqrt 2}{2}$~: $2y \in
\intcc{-\frac\pi2}{\frac\pi2}$, donc $\arcsin(\sin 2y) = 2y$ et
$h = 0$.

Pour $x \in \intoc{\frac{\sqrt 2}{2}}{1}$~: $2y \in
\intoc{\frac\pi2}{\pi}$, et l'angle de
$\intcc{-\frac\pi2}{\frac\pi2}$ dont le sinus vaut $\sin 2y$ est
$\pi - 2y$~: donc $h(x) = \pi - 4\arcsin x$. (Vérification par la
dérivée~: on y a
$h'(x) = \frac{2(1 - 2x^2)}{\abs{1 - 2x^2}\sqrt{1 - x^2}} -
\frac{2}{\sqrt{1-x^2}} = \frac{-4}{\sqrt{1 - x^2}}$, la dérivée
de $-4\arcsin x$~; et en $x = 1$, $h(1) = \arcsin 0 - 2\cdot\frac\pi2
= -\pi = \pi - 4\cdot\frac\pi2$.)

Pour $x \in \intco{-1}{-\frac{\sqrt 2}{2}}$, par imparité de $h$~:
$h(x) = -\pi - 4\arcsin x$.
\end{solution}

\begin{solution}{exo:b1:functions:11}
$g = \arctan \circ \sinh$ est une composée de fonctions impaires
strictement croissantes~: elle est donc impaire et strictement
croissante~; quand $x \to +\infty$, $\sinh x \to +\infty$ donc
$g(x) \to \frac\pi2$, et $g$ est une \omterm{def:b1:logic:inj}{bijection} sur
$\intoo{-\frac\pi2}{\frac\pi2}$ (continuité et théorème des valeurs
intermédiaires). Dérivation des fonctions composées, avec la
\cref{prop:b1:functions:hyprules}~(1)~:
\[
g'(x) = \frac{\cosh x}{1 + \sinh^2 x} = \frac{\cosh x}{\cosh^2 x}
= \frac{1}{\cosh x} .
\]
Par construction, $\tan g(x) = \sinh x$. Ensuite, comme $g(x) \in
\intoo{-\frac\pi2}{\frac\pi2}$ a un cosinus positif,
\begin{align*}
\cos g(x) &= \frac{1}{\sqrt{1 + \tan^2 g(x)}}
= \frac{1}{\sqrt{1 + \sinh^2 x}} = \frac{1}{\cosh x},\\
\sin g(x) &= \tan g(x) \cos g(x) = \frac{\sinh x}{\cosh x} = \tanh x .
\end{align*}
\end{solution}

\begin{solution}{exo:b1:functions:12}
Posons $u = \arctan 2$ et $v = \arctan 3$~; tous deux appartiennent à
$\intoo{\frac\pi4}{\frac\pi2}$ (leurs arguments dépassent $1$), donc
$u + v \in \intoo{\frac\pi2}{\pi}$. La formule d'addition donne
\[
\tan(u + v) = \frac{2 + 3}{1 - 6} = -1 ,
\]
et l'unique angle de $\intoo{\frac\pi2}{\pi}$ dont la tangente vaut
$-1$ est $\frac{3\pi}4$~: donc $\arctan 2 + \arctan 3 =
\frac{3\pi}4$. (Appliquer aveuglément $\arctan$ à la tangente
donnerait $-\frac\pi4$, faux de $\pi$ --- c'est la localisation qui
sauve le calcul, et le dénominateur négatif $1 - 6$ est précisément
le signal que la somme a quitté
$\intoo{-\frac\pi2}{\frac\pi2}$.) En ajoutant $\arctan 1 =
\frac\pi4$~:
\[
\arctan 1 + \arctan 2 + \arctan 3 = \frac\pi4 + \frac{3\pi}4 = \pi .
\]
\end{solution}

\begin{solution}{pb:b1:functions:1}
\textbf{1.} $f$ est un quotient de fonctions dérivables dont le
dénominateur ne s'annule pas sur $\intoo0{+\infty}$, et
\[
f'(t) = \frac{\frac1t \cdot t - \ln t}{t^2} = \frac{1 - \ln
t}{t^2} ,
\]
positive pour $t < \eu$, nulle en $t = \eu$, négative pour $t > \eu$~:
$f$ croît strictement sur $\intoc0\eu$, décroît strictement sur
$\intco\eu{+\infty}$, avec pour maximum $f(\eu) = \frac1\eu$.

\textbf{2.} Quand $t \to 0^+$~: $\ln t \to -\infty$ et $\frac1t \to
+\infty$, donc $f(t) \to -\infty$. Quand $t \to +\infty$~: $f(t) \to
0$ par les \omterm{prop:b1:functions:powerrules}{croissances comparées} de la
\cref{prop:b1:functions:powerrules} ($\beta = \alpha = 1$). Signe~:
celui de $\ln t$, donc $f < 0$ sur $\intoo01$, $f(1) = 0$, $f > 0$ sur
$\intoo1{+\infty}$. Le graphe monte depuis $-\infty$, traverse zéro en
$1$, culmine en $(\eu, \frac1\eu)$, puis décroît vers $0^+$.

\textbf{3.} $f(4) = \frac{\ln 4}4 = \frac{2\ln 2}4 = \frac{\ln 2}2
= f(2)$.

\textbf{4.} Par le tableau de variations et le théorème des valeurs
intermédiaires (utilisé au niveau du lycée~; formalisé au
\cref{ch:b1:continuity}). Pour $c < 0$~: les solutions n'existent que
là où $f < 0$, c'est-à-dire dans $\intoo01$, où $f$ est une \omterm{def:b1:logic:inj}{bijection}
strictement croissante sur $\intoo{-\infty}0$~: exactement une. Pour
$c = 0$~: seulement $t = 1$. Pour $0 < c < \frac1\eu$~: sur
$\intoo1\eu$, $f$ croît de $0$ à $\frac1\eu$~: une solution~; sur
$\intoo\eu{+\infty}$, $f$ décroît de $\frac1\eu$ à $0$~: une de plus~;
en tout, deux. Pour $c = \frac1\eu$~: seulement le point maximum
$t = \eu$. Pour $c > \frac1\eu$~: aucune.

\textbf{5.} Pour $x \neq \eu$, la stricte maximalité donne $f(x) <
f(\eu)$, c'est-à-dire $\frac{\ln x}x < \frac1\eu$, c'est-à-dire
$\eu\ln x < x$, c'est-à-dire $\ln(x^\eu) < \ln(\eu^x)$~: $x^\eu <
\eu^x$. Avec $x = \pi$~: $\eu^\pi > \pi^\eu$ (numériquement,
$23.14 > 22.46$).

\textbf{6.} Pour $x, y > 0$, les deux membres sont positifs, donc
\[
x^y = y^x \iff y\ln x = x\ln y \iff \frac{\ln x}x = \frac{\ln y}y
\iff f(x) = f(y) ,
\]
en divisant par $xy > 0$.

\textbf{7.} Soit $f(x) = f(y) = c$ avec $x \neq y$. Si $c \leq 0$ ou
$c = \frac1\eu$, la question 4 dit que l'équation $f = c$ a une unique
solution~: impossible. Donc $0 < c < \frac1\eu$, et, toujours par la
question 4, les deux solutions sont un point de $\intoo1\eu$ et un
point de $\intoo\eu{+\infty}$~: les deux coordonnées dépassent $1$ et
elles encadrent $\eu$.

\textbf{8.} En substituant $y = tx$ dans $y\ln x = x\ln y$~:
\[
tx\ln x = x(\ln t + \ln x)
\iff (t - 1)\ln x = \ln t
\iff \ln x = \frac{\ln t}{t - 1} ,
\]
donc $x = t^{1/(t-1)}$ et $y = tx = t^{1 + \frac1{t-1}} =
t^{t/(t-1)}$. Réciproquement, pour ces valeurs, $\ln y = \frac{t\ln
t}{t-1} = t\ln x$ et $y = tx$, donc $\frac{\ln y}y = \frac{t\ln
x}{tx} = \frac{\ln x}x$~: c'est une solution. Toute solution non
triviale a un rapport $t = y/x \in \intoo0{+\infty} \setminus \{1\}$,
donc le paramétrage est complet.

\textbf{9.} $t = 2$~: $x = 2^{1/1} = 2$, $y = 2^{2/1} = 4$. Et
\[
x(1/t) = (1/t)^{\frac1{\frac1t - 1}}
= (t^{-1})^{\frac{t}{1 - t}}
= t^{\frac{t}{t - 1}} = y(t) ,
\]
puis $y(1/t) = \frac1t\,x(1/t) = \frac{y(t)}t = x(t)$~: le changement
de paramètre $t \mapsto 1/t$ échange les coordonnées, comme l'exige
la symétrie de l'équation.

\textbf{10.} Quand $t \to 1$~: $\frac{\ln t}{t-1} \to 1$ (c'est le
taux d'accroissement de $\ln$ en $1$), donc $x \to \eu^1 = \eu$ et
$y = tx \to \eu$. Quand $t \to +\infty$~: $\ln x = \frac{\ln t}{t-1}
\to 0$ donc $x \to 1$, tandis que $\ln y = \frac{t\ln t}{t-1} \to
+\infty$ donc $y \to +\infty$. Quand $t \to 0^+$~: $\ln x = \frac{\ln
t}{t-1} \to \frac{-\infty}{-1} = +\infty$ donc $x \to +\infty$,
tandis que $\ln y = \frac{t\ln t}{t-1} \to \frac{0}{-1} = 0$ donc $y
\to 1$ (en utilisant $t\ln t \to 0$, \cref{ex:b1:functions:xx}). La
branche court donc depuis l'asymptote $y = 1$ (tout à droite), monte
en traversant $(\eu, \eu)$ sur la diagonale, et s'échappe le long de
l'asymptote $x = 1$ (tout en haut) --- symétrique par rapport à la
diagonale d'après la question 9.

\textbf{11.} $\ln x(t) = g(t) = \frac{\ln t}{t-1}$, et
\[
g'(t) = \frac{\frac{t-1}t - \ln t}{(t-1)^2}
= \frac{h(t)}{(t-1)^2},
\qquad h(t) = 1 - \frac1t - \ln t .
\]
$h(1) = 0$ et $h'(t) = \frac1{t^2} - \frac1t = \frac{1 -
t}{t^2} < 0$ pour $t > 1$~: $h < 0$ sur $\intoo1{+\infty}$, donc
$g' < 0$ et $x = \eu^g$ y est strictement décroissante (de $\eu$ à
$1$, d'après la question 10).

\textbf{12.} D'après la question 11, $t \mapsto x(t)$ est une
\omterm{def:b1:logic:inj}{bijection} strictement décroissante de $\intoo1{+\infty}$ sur
$\intoo1\eu$~; notons $t(x)$ sa réciproque (également strictement
décroissante) et posons $\varphi(x) = y(t(x))$. Remarquons que
$h < 0$ également sur $\intoo01$ (on y a $h' > 0$ et $h(1) = 0$),
donc $x(\cdot)$ est aussi décroissante sur $\intoo01$~; par
conséquent $y(t) = x(1/t)$ (question 9) est \emph{croissante} en $t$
sur $\intoo1{+\infty}$, de $\eu$ à $+\infty$. Par composition,
$\varphi = y \circ t(\cdot)$ est strictement décroissante de
$\intoo1\eu$ sur $\intoo\eu{+\infty}$, et $x^{\varphi(x)} =
\varphi(x)^x$ d'après la question 8. Enfin, pour une valeur commune
$c = f(x) \in \intoo0{\frac1\eu}$, la paire $\{x, \varphi(x)\}$ est
\emph{l'}\omterm{def:b1:logic:sets}{ensemble} des deux solutions de $f = c$ (question 4)~; en
prolongeant $\varphi$ à $\intoo\eu{+\infty}$ par l'\omterm{def:b1:logic:map}{application}
réciproque, $\varphi(\varphi(x))$ redonne l'autre élément
(c'est-à-dire l'élément initial)~: $\varphi \circ \varphi =
\mathrm{id}$.

\textbf{13.} Si $(x, y)$ est une solution entière non triviale avec
$x < y$, la question 7 place $x \in \intoo1\eu$~: le seul entier de
cet intervalle est $x = 2$. Alors $f(y) = f(2) = \frac{\ln2}2$ avec
$y > \eu$~; d'après la question 4, l'équation $f = \frac{\ln2}2$ a
exactement une solution au-delà de $\eu$, et la question 3 l'exhibe~:
$y = 4$. Donc $(2, 4)$, et son échangé, sont les seules.

\textbf{14.} $t = 1 + \frac1n = \frac{n+1}n$ donne $\frac1{t - 1}
= n$ et $\frac t{t-1} = n + 1$, donc
\[
x_n = \Bigl(\frac{n+1}n\Bigr)^{n},
\qquad
y_n = \Bigl(\frac{n+1}n\Bigr)^{n+1},
\]
manifestement rationnels, distincts ($y_n = t\,x_n \neq x_n$), et
$n = 1$ donne $x_1 = 2$, $y_1 = 4$.

\textbf{15.} $t = y/x$ est rationnel et $> 1$~; écrivons $t - 1 =
\frac rs$ sous forme irréductible, de sorte que $t = \frac{s + r}s$
avec $\gcd(s + r, s) = \gcd(r, s) = 1$. La question 8 donne
$x = t^{s/r}$, d'où
\[
x^r = t^s = \frac{(s+r)^s}{s^s} ,
\]
fraction irréductible (aucun nombre premier ne divise à la fois
$s + r$ et $s$). En écrivant $x = \frac pq$ sous forme irréductible,
$x^r = \frac{p^r}{q^r}$ l'est aussi, et, par unicité de la
représentation irréductible~: $p^r = (s+r)^s$ et $q^r = s^s$.
Comparons l'exposant d'un nombre premier $\ell$ quelconque dans
$q^r = s^s$~: $r \cdot v_\ell(q) = s \cdot v_\ell(s)$, donc $r$
divise $s\,v_\ell(s)$~; comme $\gcd(r, s) = 1$, $r$ divise
$v_\ell(s)$ pour tout nombre premier $\ell$ (décomposition en
facteurs premiers unique, admise~; démontrée au \cref{ch:b1:arith}),
donc $s = b^r$ pour un entier $b$. Le même argument sur
$p^r = (s+r)^s$ donne $s + r = a^r$.

\textbf{16.} Supposons $a^r - b^r = r$ avec des entiers $a > b \geq 1$
et $r \geq 2$. Alors $a \geq b + 1$, donc, par la formule du binôme,
\[
a^r - b^r \geq (b+1)^r - b^r
= \sum_{k=0}^{r-1}\binom rk b^k
\geq \binom r0 + \binom r{r-1} b^{r-1}
\geq 1 + r > r ,
\]
ce qui est contradictoire. D'après la question 15, $a^r - b^r =
(s + r) - s = r$ impose $r = 1$~: $t = 1 + \frac1s$, et la solution
est $(x_s, y_s)$ de la question 14. Avec les couples échangés, la
classification est complète.

\textbf{17.} $f\bigl(\tfrac94\bigr) = \frac{\ln 2.25}{2.25} =
\frac{0.81093}{2.25} = 0.3604$ et $f\bigl(\tfrac{27}8\bigr) =
\frac{\ln 3.375}{3.375} = \frac{1.21640}{3.375} = 0.3604$ (à quatre
décimales)~: égaux, comme la construction le promettait --- donc
$\bigl(\tfrac94\bigr)^{27/8} = \bigl(\tfrac{27}8\bigr)^{9/4}$.

\textbf{18.} Les paramètres $t_n = 1 + \frac1n$ décroissent
strictement vers $1$. Comme $x(\cdot)$ est strictement décroissante
sur $\intoo1{+\infty}$ (question 11), $x_n = x(t_n)$ est strictement
\emph{croissante}~; comme $y(\cdot)$ y est strictement croissante
(question 12), $y_n = y(t_n)$ est strictement \emph{décroissante}.
D'après la question 7, $x_n \in \intoo1\eu$ et $y_n \in
\intoo\eu{+\infty}$~: donc $x_n < \eu < y_n$ pour tout $n$. Enfin,
$t_n \to 1$ et la question 10 donnent $x_n \to \eu$ et $y_n \to \eu$.
La convergence monotone classique de $\bigl(1 + \frac1n\bigr)^n$
tombe de la géométrie de la branche, sans aucune inégalité nouvelle.

\textbf{19.} Si $\eu \leq a < b$~: $f$ strictement décroissante sur
$\intco\eu{+\infty}$ donne $f(a) > f(b)$, c'est-à-dire $b\ln a > a\ln
b$, c'est-à-dire $a^b > b^a$. Si $1 < a < b \leq \eu$~: $f$
strictement croissante donne $f(a) < f(b)$, d'où $a^b < b^a$. Cas
mixte~: $2 < \eu < 3$ et $2^3 = 8 < 9 = 3^2$~; mais $2 < \eu < 5$ et
$2^5 = 32 > 25 = 5^2$~: les deux issues se produisent, selon le côté
de la branche où tombe le point $(a, b)$.

\textbf{20.} La branche rencontre la diagonale en $\eu$ et $\varphi$
s'y prolonge continûment avec $\varphi(\eu) = \eu$. En dérivant
$\varphi(\varphi(x)) = x$ en $x = \eu$ par la règle de dérivation des
fonctions composées~: $\varphi'(\varphi(\eu))\,\varphi'(\eu) =
\varphi'(\eu)^2 = 1$, donc $\varphi'(\eu) = \pm1$~; $\varphi$ est
décroissante, donc $\varphi'(\eu) = -1$. La branche coupe la
diagonale avec la pente $-1$~: perpendiculairement.

\textbf{21.} $y = 3x$ est le cas $t = 3$~: $x = 3^{1/2} = \sqrt3$
et $y = 3^{3/2} = 3\sqrt3$. Vérification~: $f(\sqrt3) =
\frac{0.54931}{1.73205} = 0.3171$ et $f(3\sqrt3) =
\frac{\ln 5.19615}{5.19615} = \frac{1.64792}{5.19615} = 0.3171$~:
égaux à quatre décimales, donc $(\sqrt3)^{3\sqrt3} =
(3\sqrt3)^{\sqrt3}$.

\textbf{22.} $n^{1/n} = \eu^{f(n)}$, et $\exp$ est croissante, donc
l'ordre est celui des $f(n)$~: $f(3) = \frac{\ln3}3 \approx
0.3662$ l'emporte sur $f(2) = f(4) = \frac{\ln2}2 \approx 0.3466$
(question 3) et sur $f(5) \approx 0.3219$. Le plus grand est
$\sqrt[3]3$, et le couple égal est $\sqrt2 = \sqrt[4]4$ --- le couple
entier $(2, 4)$ sous un déguisement de plus.

\textbf{23.} L'\omterm{def:b1:logic:sets}{ensemble} des solutions est la réunion de la diagonale
$\{(x, x) : x > 0\}$ et d'une unique branche, symétrique par rapport
à la diagonale, strictement décroissante depuis l'asymptote $x = 1$
(quand $y \to +\infty$) jusqu'à l'asymptote $y = 1$ (quand $x \to
+\infty$), coupant la diagonale exactement une fois, en $(\eu, \eu)$,
avec la pente $-1$ en ce point. Sur la branche se trouvent les points
entiers $(2, 4)$ et $(4, 2)$ --- les seuls --- et les points
rationnels $(x_n, y_n)$, qui progressent monotonement le long de la
branche vers $(\eu, \eu)$ sans jamais l'atteindre ($\eu$ est
irrationnel~; les points rationnels de la branche s'accumulent en un
coin irrationnel).

\textbf{24.} (i) Les \omterm{prop:b1:functions:powerrules}{croissances comparées} ont donné les limites de
$f$ en $+\infty$ (question 2) et $t\ln t \to 0$ (question 10),
façonnant à la fois le tableau de variations et les asymptotes. (ii)
Le théorème des valeurs intermédiaires, à travers les énoncés de
\omterm{def:b1:logic:inj}{bijectivité}, a converti le tableau de variations en comptages exacts
de solutions (question 4) et en existence de l'\omterm{def:b1:logic:map}{application} réciproque
$t(x)$ (question 12). (iii) L'unicité de la décomposition en facteurs
premiers a alimenté les deux identifications de fractions
irréductibles de la question 15, cœur arithmétique de la
classification des points rationnels.

\textbf{25.} Un seul calcul de dérivée --- le signe de $1 - \ln
t$ --- a tout engendré~: le nombre de solutions pour chaque niveau
$c$, la forme et les asymptotes de la branche, l'inégalité extrémale
$\eu^x \geq x^\eu$, la classification des solutions entières et
rationnelles, et la convergence monotone de $\bigl(1 +
\frac1n\bigr)^n$. Telle est l'économie de la pensée par tableaux de
variations~: une étude en dimension un résout une équation à deux
variables, parce que l'équation se factorise à travers une unique
fonction. Le \cref{ch:b1:seq} refera la convergence de $(x_n)$ avec
la théorie générale des suites monotones~; le \cref{ch:b1:derivative}
fonde rigoureusement les tableaux de variations sur le théorème des
accroissements finis~; et le \cref{ch:b1:taylor} fournit ce que le
tableau ne peut donner --- la \emph{vitesse} de la convergence
$x_n \to \eu$ (elle est d'ordre $1/n$).
\end{solution}
